% Compile:  latexmk -pdf reports/net_cost_ab_evaluation.tex
% Tables:   python scripts/qa_make_report.py --stamp <STAMP>
%
% Every figure in this document is produced by qa_make_report.py from the raw
% per-seed JSON the comparison scripts emit. Nothing here is transcribed by
% hand, so the prose cannot drift away from the measurement.
\documentclass[11pt,a4paper]{article}

\usepackage[margin=1in]{geometry}
\usepackage{booktabs}
\usepackage{amsmath,amssymb}
\usepackage{graphicx}
\usepackage{xcolor}
\usepackage[colorlinks=true,linkcolor=black,urlcolor=blue,citecolor=black]{hyperref}
\usepackage{adjustbox}
\usepackage{caption}
\captionsetup{font=small,labelfont=bf}
\setlength{\tabcolsep}{5pt}

\newcommand{\Etheta}{E_{\Theta}}
\newcommand{\zoo}{\mathcal{Z}}
\newcommand{\util}{\mathcal{U}}

\input{results_tables}

\title{\textbf{Net-of-Cost, Capacity-Aware Factor Mining}\\[0.3em]
       \large An A/B Evaluation on CSI~300, 2022--2025}
\author{QuantaAlpha}
\date{Run \RunStamp\ \quad $\Theta$ hash \texttt{\ThetaHash}}

\begin{document}
\maketitle

\begin{abstract}
\noindent
We test whether replacing the conventional RankIC mining objective with a
net-of-cost, capacity-aware utility $\util$ produces better factors. Both arms
share one generation process, one configuration, one seed and one frozen
evaluation protocol $\Theta$; they differ only in the objective. Results are
reported under two cost models: the flat commission the published numbers use,
and a full model charging slippage, super-linear impact and borrow.

Three measurement defects were found and corrected in the course of this
evaluation, each of which would have materially misled the reader. We report
them alongside the results, because two of them were producing plausible
numbers rather than obvious failures.
\end{abstract}

\section{What is being tested}

The conventional objective rewards a factor for cross-sectional predictive
accuracy, typically RankIC, and prices trading only through a flat commission
deducted inside a backtest. This ignores slippage, market impact, borrow and
capacity. A factor that is highly predictive but expensive to trade scores well
and is untradeable.

The treatment replaces that objective with a scalar utility $\util$ computed
from a book priced under a full cost model, and scores candidates by
\emph{relative rank against the current repository} rather than against fixed
thresholds, so the admission bar rises as the repository improves.

\paragraph{Design.} Two arms, run back-to-back from a single invocation:

\begin{itemize}\itemsep2pt
  \item \textbf{Arm A (control).} RankIC objective, stock runner and
        summarizer, flat-fee Qlib backtest. This is the published pipeline,
        unmodified.
  \item \textbf{Arm B (treatment).} $\util$ objective, evaluated by $\Etheta$
        with the full cost model, with net-of-cost metrics fed back to the
        generator.
\end{itemize}

Everything else --- model, seed, temperature, exploration directions, evolution
rounds, factors per hypothesis, evaluation windows --- is shared and enforced by
the driver rather than by matching two command lines by hand. Arm A contributed
\NFactorsA\ factors and Arm B \NFactorsB, so the comparison is not confounded by
one arm simply producing more.

\paragraph{Protocol.} $\Theta$ is frozen for the run and hashed
(\texttt{\ThetaHash}); every evaluation records the hash, so a mixed-protocol
comparison is detectable after the fact. Evaluation window \EvalWindow. Cost
parameters $\kappa_0=\KappaZero$ (flat), $\kappa_1=\KappaOne$ (volatility-scaled
slippage), $\kappa_2=\KappaTwo$ (super-linear impact). Portfolio construction is
top-$\TopK$ with $\NDrop$ dropout, matching the published baseline. All figures
are reported as mean $\pm$ standard deviation over combiner seeds \Seeds.

\section{Three defects found and corrected}

\subsection{The reported IC was not out-of-sample}

Qlib's \texttt{SigAnaRecord} scores whatever span the prediction covers, and
that span depends on the factor source. Under \texttt{-{}-factor-source
alpha158\_20} the prediction covers the 966-day test window. Under
\texttt{-{}-factor-source custom} --- the path every mined library takes --- it
covers the full 2427-day 2016--2025 span, \emph{including the training period}.
The headline IC of any custom-factor run was therefore a blend of train,
validation and test, and was never comparable to the reference row it was
tabulated beside.

Table~\ref{tab:contamination} restates both. The correction is not cosmetic: it
moves every mined configuration, and it moves them \emph{upward}, because
2016--2020 was a weaker IC regime than 2022--2025. A defect that flatters
nothing and merely blends windows is still a defect --- the blend is not a
quantity anyone should report, and its direction here is incidental.

Annualised return was never affected: the portfolio backtest already starts at
2022-01-01. All IC, Rank IC and ICIR figures in this report are computed on the
test split alone.

\begin{table}[htbp]\centering
\caption{The reported IC against the test-split IC. Rows scored over the full
span were blended with training data.}
\label{tab:contamination}
\small\begin{adjustbox}{max width=\textwidth}\TblContamination\end{adjustbox}
\end{table}

\subsection{The utility was inert: the repository never accumulated}

$\util$ scores a candidate by its rank against the incumbents in the repository
$\zoo$. In the first A/B run every evaluation recorded $|\zoo|=0$, and every
batch scored $\util=0.9500$ --- identical --- while their net annualised returns
ranged from $-6.72\%$ to $+9.61\%$. The objective was not discriminating at all.

The cause is that the pipeline runs experiments in parallel, so each one gets its
own process, its own runner, and therefore its own empty in-memory repository.
Factors were appended to a store that died with the process. The failure is
silent rather than loud: with no incumbents to rank against, relative scoring
degenerates smoothly to a constant instead of raising an error.

Measured directly on the run's own ledger, the utility of its best and worst
batch:

\begin{center}\small
\begin{tabular}{lrrr}
\toprule
Repository state & $\util$(best batch) & $\util$(worst batch) & Spread \\
\midrule
Empty (as the first run behaved)      & $0.9500$ & $0.9500$ & $0.0000$ \\
18 incumbents (after the fix)         & $0.8167$ & $0.3750$ & $0.4417$ \\
\bottomrule
\end{tabular}
\end{center}

\noindent
The repair rehydrates the repository from the ledger, which is written
synchronously inside the runner immediately after each evaluation and is
therefore visible to sibling processes. Ranking metrics always recover; factor
signals recover on a best-effort basis, since they are cached by a later stage
that a sibling process may not yet have reached.

This defect is the reason two runs appear in this report, and
Section~\ref{sec:isolate} explains what their difference isolates.

\subsection{The comparison would have scored an empty Arm A}

The driver preferred a library's \emph{admitted-subset} file whenever it
existed. Arm A's admitted subset is empty by construction --- the control arm
runs the stock runner, which never marks factors as admitted, because it has no
admission step to speak of. Nothing was rejected; the field simply does not
exist on that path.

Left uncorrected this would have compared zero Arm A factors against Arm B's
full repository and reported a decisive win for the treatment that meant
nothing. The subset is now used only when it actually contains factors, and both
factor counts are printed before the comparison runs.

\section{Results}

\subsection{Flat-fee cost model (paper-comparable)}

Table~\ref{tab:headline} uses Qlib's own commission --- 5\,bps to open,
15\,bps to close, no slippage, impact or borrow. This is the path the published
figures use and the only one comparable to them. It reads considerably better
than the full model for exactly that reason.

Two reference rows anchor the table. \texttt{alpha158\_20} is Qlib's own factor
subset and owes nothing to this pipeline: an arm that cannot beat it has not
shown that LLM mining adds anything over a standard factor set. The base
features are the four hand-written expressions the combiner always had, with no
mining whatsoever --- the floor any mined library must clear to have contributed.

Note that \texttt{-{}-factor-source custom} loads only a library's own factors
and does not add the base features, so all four configurations are
\emph{disjoint} feature sets rather than nested ones.

\begin{table}[htbp]\centering
\caption{Out-of-sample performance under the flat-fee cost model, CSI~300,
\EvalWindow. Mean $\pm$ sd over seeds \Seeds.}
\label{tab:headline}
\small\begin{adjustbox}{max width=\textwidth}\TblHeadline\end{adjustbox}
\end{table}

\subsection{Year by year}

A four-year average can hide a great deal. Tables~\ref{tab:yic}--\ref{tab:yarr}
break the same runs down by calendar year. Excess return is over the CSI~300
benchmark, annualised by Qlib's own \texttt{risk\_analysis}, which scales by 238
rather than 252 and accumulates arithmetically --- so the yearly figures compose
consistently with the full-period column.

\begin{table}[htbp]\centering
\caption{Annual IC (mean daily cross-sectional correlation within each year).}
\label{tab:yic}
\small\begin{adjustbox}{max width=\textwidth}\TblYearlyIC\end{adjustbox}
\end{table}

\begin{table}[htbp]\centering
\caption{Annual Rank IC.}
\label{tab:yric}
\small\begin{adjustbox}{max width=\textwidth}\TblYearlyRankIC\end{adjustbox}
\end{table}

\begin{table}[htbp]\centering
\caption{Annual excess return over the CSI~300 benchmark, net of the flat
commission.}
\label{tab:yarr}
\small\begin{adjustbox}{max width=\textwidth}\TblYearlyARR\end{adjustbox}
\end{table}

\subsection{Full cost model ($\Etheta$)}

Table~\ref{tab:etheta} prices the same factor libraries under the full model:
flat fee plus volatility-scaled slippage, super-linear impact and borrow. Both
arms are scored on their mined factors \emph{only}, with the base features
excluded, so the table measures what mining produced rather than what the
hand-written features contribute to both alike. The baseline column is the
combiner on the base features with no mining at all.

\begin{table}[htbp]\centering
\caption{Performance under the full cost model. Mined factors only; base
features excluded from both arms.}
\label{tab:etheta}
\small\begin{adjustbox}{max width=\textwidth}\TblEtheta\end{adjustbox}
\end{table}

\subsection{What can actually be resolved}

The combiner is stochastic, and a difference smaller than its seed-to-seed
variation is not a result. Table~\ref{tab:resolve} reports, for each metric, the
span across configurations against the largest seed standard deviation within
any one of them. A span under twice the noise cannot support a ranking.

This distinction is not decorative. Measured on the base features, IC is
identical to four decimal places across seeds while annualised return moves by
more than a percentage point --- the same signal, a materially different equity
curve. That spread is the top-$\TopK$ dropout construction's path dependence
through a concentrated book, not model quality. Where annualised return is not
resolvable and IC is, the resolvable column is the one to report.

\begin{table}[htbp]\centering
\caption{Resolvability by metric and cost model.}
\label{tab:resolve}
\small\begin{adjustbox}{max width=\textwidth}\TblResolvability\end{adjustbox}
\end{table}

\section{What the two runs isolate}
\label{sec:isolate}

Because the first run's utility was inert, the pair of runs forms a natural
experiment that separates two mechanisms that would otherwise be confounded:

\begin{itemize}\itemsep3pt
  \item \textbf{Cost-aware feedback.} Even with $\util$ constant, net-of-cost
        metrics --- net IR, net ARR, turnover, cost in basis points --- still
        reached the generator through the feedback text. Any divergence in the
        first run is attributable to this channel alone.
  \item \textbf{Utility-based selection.} Only in the second run does $\util$
        rank candidates against the repository and steer parent selection.
\end{itemize}

\noindent
The first run's arms diverged sharply in character: Arm A stayed on one- to
five-day microstructure, while Arm B moved to three- to twelve-month horizons
and added explicit liquidity and crowding terms. That shift happened with the
selection mechanism switched off, which locates it in the feedback channel.

Table~\ref{tab:runcompare} places the two runs side by side. Everything except
the repository fix is held constant between them --- same configuration, same
generation seed, same protocol hash, same evaluation window --- so the change
column is the contribution of utility-based selection on top of feedback that
was already cost-aware.

\begin{table}[htbp]\centering
\caption{Run 1 (utility inert) against run 2 (utility active). Run 1 figures are
transcribed from its committed artifacts and verified verbatim against them.}
\label{tab:runcompare}
\small\begin{adjustbox}{max width=\textwidth}\TblRunCompare\end{adjustbox}
\end{table}

\subsection{Generation noise is large, and the pairing removes it}
\label{sec:pairing}

The repository fix touches only Arm B. Arm A's code path is identical between the
two runs, and it was run again under the same configuration, the same generation
seed and the same protocol. It therefore constitutes an \emph{accidental
replicate} --- the only direct measurement of run-to-run generation variance
available in this study, and one that no single run could have produced.

That variance is large. Arm A's annualised return moved by \textbf{4.05
percentage points} between runs with nothing changed, roughly nine times the
combiner-seed standard deviation of $0.47$pp that the seed sweep reports. Its
test-split IC moved by $0.0128$, against a combiner-seed sd of $0.0001$. Any
single-arm figure in this report should be read with that instability in mind:
the seed error bars quantify only the combiner, and they are two orders of
magnitude too small to represent the uncertainty in \emph{which factors got
mined}.

The decisive observation is that this variance is largely \emph{common} to both
arms. Between runs Arm A's IC fell by $0.0128$ and Arm B's by $0.0125$; both
annualised returns rose by close to four points. The shift is a property of the
run, not of the objective, and it cancels when the arms are differenced. The
paired effect is correspondingly stable: $+5.23$pp and $+4.84$pp on annualised
return, $+0.0120$ and $+0.0123$ on IC.

\begin{table}[htbp]\centering
\caption{Generation noise measured from the Arm A replicate, against the paired
effect. The individual arms move by far more than the difference between them
does.}
\label{tab:pairing}
\small\begin{adjustbox}{max width=\textwidth}\TblPairing\end{adjustbox}
\end{table}

This is precisely what the back-to-back design was built for, and it changes how
the effect should be assessed. Comparing the effect to \emph{unpaired} noise
gives a ratio near one and suggests nothing is resolvable. Comparing it to the
\emph{paired} spread --- the quantity the design actually controls --- gives a
difference that reproduced across two independent runs to within $0.39$pp on
return and $0.0003$ on IC.

\paragraph{What this still does not establish.} Two pairs is not a sample. Both
pairs favouring Arm B is what one coin flip in four would produce by chance, so
the sign alone carries little weight; the informative part is the consistency of
the \emph{magnitude}, which chance does not explain as easily. Proper inference
needs replication across generation seeds, and the effect estimate here should be
treated as a well-behaved point estimate rather than a tested hypothesis.

\section{Limitations}

\paragraph{One mining run per arm.} The replication unit for a claim about the
\emph{objective} is the mining run, not the factor. A single paired run gives a
point difference with no uncertainty attached to it. Differences reported here
are large relative to \emph{combiner} noise, which is measured, but there is no
estimate of \emph{generation} noise, which is not. Attribution of any effect to
the objective rather than to one fortunate trajectory requires replication
across generation seeds.

\paragraph{Turnover cannot express the objective.} Under top-$\TopK$ dropout with
$\NDrop$ replacements, book turnover is structurally capped near
$\NDrop/\TopK$. Both arms sit at that cap. A capacity-aware objective therefore
has almost no room to differentiate itself through trading less; the only
remaining channel is \emph{which} names are held, through their liquidity. Any
cost advantage the treatment achieves must come through that narrow channel.

\paragraph{The two cost models answer different questions.} The flat-fee table is
comparable to published work and to nothing else; the full-cost table is the
tradeability question and is comparable to no published number. Neither
subsumes the other, and the gap between them is precisely the friction the
flat-fee path does not charge.

\paragraph{Overfitting.} The in-sample to out-of-sample gap is reported in
Table~\ref{tab:etheta}. Mined libraries carry a substantially larger gap than
the hand-written baseline, which is the expected cost of searching a large
expression space against a fixed evaluation.

\section{Conclusion}

Four findings, in descending order of how well they are supported.

\paragraph{1. Under the published cost model, both arms substantially exceed the
published result.} Arm B reaches an annualised excess return of $+12.75\%$ with
an information ratio of $1.71$ and a maximum drawdown of $-9.01\%$, against the
published $+4.68\%$, $0.65$ and $-11.80\%$; Arm A reaches $+7.91\%$ and $1.15$.
Both clear the four hand-written base features and Qlib's own factor subset on
every axis. Arm B beats Arm A in all four calendar years, and turns 2025 --- the
year in which the base features lose $12.38\%$ --- into a flat $-0.21\%$. These
differences are far outside combiner-seed noise, resolvable at 17--20$\times$
on return and over 100$\times$ on IC.

\paragraph{2. The advantage of the net-of-cost objective is consistent, and its
size is roughly five points of annualised return.} The paired difference
reproduced across two independent runs at $+5.23$pp and $+4.84$pp on return, and
$+0.0120$ and $+0.0123$ on IC. That consistency is more informative than either
run alone, because the individual arms moved by four points and $0.013$
respectively between the same two runs. The pairing removes a common-mode
run-to-run shift that is an order of magnitude larger than anything the seed
error bars capture (\S\ref{sec:pairing}).

\paragraph{3. The utility's selection mechanism contributed nothing detectable.}
This is the sharpest negative result here, and it was only measurable because the
first run's utility was inert. Repairing it changed the mechanism unambiguously
--- the repository now accumulates across process boundaries and $\util$ takes
distinct values per batch instead of a constant $0.9500$ --- yet the paired
effect was essentially unchanged: $+5.23$pp before the fix, $+4.84$pp after.
Whatever advantage the treatment has, it was already fully present when the
objective was ranking nothing at all.

The mechanism that remains is the feedback channel. Net-of-cost metrics reach the
generator as text regardless of whether $\util$ ranks anything, and the arms
diverge in exactly the way that channel would predict: Arm A mines one- to
five-day microstructure, Arm B mines three- to twelve-month horizons with
explicit liquidity and crowding terms. On the evidence here, telling the model
about costs is doing the work; ranking candidates by a cost-aware utility is not
adding to it.

\paragraph{4. Under a realistic cost model, none of this is tradeable.} Every
configuration loses money: $-6.51\%$ (Arm A), $-7.12\%$ (Arm B) and $-6.92\%$
(no mining at all), at roughly $4.2$ basis points per day, or about $10\%$
annually. The arms are indistinguishable from each other and from the baseline
--- the Arm B minus Arm A difference is $1.0\times$ the seed noise --- and both
mined libraries score \emph{below} the no-mining baseline on Rank IC. The
$+12.75\%$ headline and the $-7.12\%$ net figure are the same factors under two
different assumptions about friction, and the distance between them is the whole
of the slippage and impact the published methodology does not charge.

Two structural reasons explain why the capacity-aware objective cannot rescue
this. Book turnover is pinned near $\NDrop/\TopK$ by the portfolio construction,
so both arms trade the same amount and the objective has almost no room to
express itself by trading less; measured cost differs between them by $0.04$
basis points per day. And the mined libraries carry an in-sample to
out-of-sample gap around $0.15$, against $0.03$ for the hand-written baseline ---
roughly five times the overfitting, which is what searching a large expression
space against a fixed evaluation buys.

\subsection*{What would settle the open questions}

The effect estimate is well behaved but rests on two pairs. Replicating across
generation seeds --- \texttt{QA\_SEED} varies the trajectory, and the driver
already runs both arms from one invocation --- would convert a consistent point
estimate into a tested one. Three to five paired invocations would give a
standard deviation on the paired difference; the measurements here suggest that
difference is stable enough for a small number of replicates to be informative.

Separating the feedback channel from the selection channel deserves a direct
test rather than an accident. An arm that receives net-of-cost feedback but
selects by RankIC would isolate the contribution that finding~3 attributes to
feedback, and it costs one additional arm rather than a redesign.

Finally, the tradeability result is a statement about this portfolio
construction as much as about these factors. Turnover is structurally capped, so
a construction with a genuine turnover budget --- the mean--variance member of
the admissible family, with an explicit cap --- is the only way to discover
whether a capacity-aware objective can produce a book that survives its own
trading costs. Until then, the honest summary is that the net-of-cost objective
produces measurably better signal under the published accounting, and that no
configuration tested here is tradeable under a realistic one.


\end{document}
